\section{第II部：仕組みから手法へ}
\label{sec:part2-mechanism-to-method}

第I部の結果，状況は次のとおりである．
中心化と切片の当てはめ直しが判断基準を扱い，ソースは向きについておよそ10本のラベル分を補い，残る向きの誤差が主な損失である．
第II部では，ソースから作れる最も強い出発点を構成し，少数のラベルも，どの発話にラベルを付けるかをラベルを使わずに選ぶことも，その向きをあまり動かせないことを示す．
そのうえで，別のモダリティを見る先生がラベルのないプール全体に付ける弱いラベルによって，足りない向きの情報を補う（図\ref{fig:method}）．

\begin{figure*}[t]
\centering
\FIG{method}{\textwidth}
\caption{提案法．ソースとターゲットの特徴は別々に中心化する．ターゲットのコーパスの外のすべてのソースから出発点の分類器$(W_0, b_0)$を作り，生徒はそこから出発し，重みと切片に別々の$\ell_2$アンカーをもつ．基本の手法では，ソースだけで学習した先生（単一モダリティの生徒に対しては，生徒自身のモダリティと生徒がもたないモダリティを含む，その構成の融合の出発点）が，ラベルのないターゲットの発話すべてに，感情ごとの合計をそろえたソフトラベルを付ける．この弱いラベルは，$4K$本の本当のラベルと並んで，合計の重み$m$で生徒の損失に入る．出発点の5つの値（アンカー$C'$と$C_b$を含む）は仮のターゲットで1度だけ選び，先生と$m$は$K$ごとに選ぶ．拡張版は，先生の候補として，$K$本のラベルで学び直した融合の生徒と，プールで学習した単一モダリティの生徒の積を加え，弱いラベルの重みづけ（釣り合わせる，釣り合わせない，ハード）を$m$とともに$K$ごとに選ぶ．}
\label{fig:method}
\end{figure*}

\subsection{出発点：中心化とアンカーつきの転移}
\label{sec:part2-start}

生徒は中心化した特徴を使う．
各ソースはそれ自身の平均で，ターゲットはそのラベルのないプールの平均で中心化し，これによって判断基準による損失の大部分が取り除かれる（\ref{sec:part1-fix}節）．
出発点$(W_0, b_0)$は，ターゲットのコーパスの外のすべてのソースの向きを合わせたものである．
各ソース$s$は，閉じた形の判別の向き$W_s(\rho_s) = ((1-\rho_s)\Sigma_s + \rho_s \tau_s I)^{-1}\tilde{M}_s$を与える．
ここで$\tilde{M}_s$はその中心化したクラス平均を並べた行列，$\Sigma_s$はその共分散行列，$\tau_s$はその固有値の平均である．
これらの向きを$N_s^{\gamma}$に比例する重みで足し合わせ，切片はクラスの事前確率が等しいと仮定して決める．
生徒は次の問題を解く．
\begin{equation}
\label{eq:student}
\min_{W, b}\ \sum_{i \in \mathcal{L}} \ell\!\left(W \mathbf{z}_i + b, y_i\right) + \frac{\|W - W_0\|_{M}^2}{2C'} + \frac{\|b - b_0\|^2}{2C_b}
\end{equation}
ここで$\mathcal{L}$は$4K$本のラベル付き発話の集合であり，ノルム$\|\cdot\|_M$は，ラベルのないターゲットのプールの散布行列から混合の重み$\rho_M$で作った計量$M$を用いる．
重みと切片に別々のアンカーを置くのは第I部の結果による．
少数のラベルで切片はよく決まるので切片には弱いアンカーで足り，少数のラベルでは重みがうまく決まらないので重みには強いアンカーが要る．
5つの値$(\gamma, \rho_s, C_b, C', \rho_M)$は，それぞれ3つの候補から仮のターゲットで選ぶ（\ref{sec:setting-selection}節）．

この構成の各段階はUARを上げる．
中心化を伴う，合併したソースからの仮説転移は，音声の$K \in \{1, 5, 10\}$で，ターゲットのみの学習，合併，中心化した合併，ターゲットの重みを選んだ合併のうち最良のものを\NUM{p2-e31-p4-vs-best-audio-lo}--\NUM{p2-e31-p4-vs-best-audio-hi}ポイント上回る（30回の反復）．
5つの値を選ぶと，$K=5$でA，A（HuBERT），V，AVにわたって\NUM{p2-e39-vs-p4-k5-lo}--\NUM{p2-e39-vs-p4-k5-hi}ポイント上がり，ソースを等しく重みづけした出発点に対しては\NUM{p2-e39-vs-d1-k5-lo}--\NUM{p2-e39-vs-d1-k5-hi}ポイント上がる（30回の反復）．
この結果を，5つの値を選んだ出発点（five-value start）と呼ぶ．
第II部と\ref{sec:comparison}節では，すべての追加をこの出発点を基準に測る．

\subsection{向きの情報を得られないところ}
\label{sec:part2-limits}

$W^{*}$をターゲットのプールのすべてのラベルを使ったときに得られる生徒とし，$E_0 = W^{*} - W_0$を出発点の向きの誤差とする．
$K$本のラベルが$E_0$に対してできることは，2つの事実によって制限される．

\begin{proposition}[張る空間]
\label{prop:span}
式\eqref{eq:student}の解の$W - W_0$の各行は，$4K$個のラベル付きの点$M^{-1}\mathbf{z}_i$（$i \in \mathcal{L}$）が張る空間に含まれる．
\end{proposition}

損失は積$W\mathbf{z}_i$を通してだけ$W$に依存するので，この命題は，$W - W_0$に表現定理\cite{kimeldorf1971results,scholkopf2001generalized}を当てはめた直接の帰結である．
修正できるのは，$E_0$のうち高々$4K$次元の部分空間への射影だけであり，$E_0$のどの部分がその部分空間に入るかは，どの発話がラベルをもつかによる．

\begin{proposition}[ラベルを使わない計画]
\label{prop:design}
損失を$W_0$のまわりのガウスの尤度で，アンカーを事前分布$W \sim \mathcal{N}(W_0, C' M^{-1})$で近似する．
この事前分布の下で$E_0$に特に好まれる向きが無いならば，任意の$4K$個の点が張る空間が捉える$\|E_0\|^2$の割合の期待値は，点のラベルを使わずに選んだどの選び方でも同じである．
\end{proposition}

ガウスの線形モデルでは事後共分散が観測されたラベルによらないので，最良の計画もラベルによらない\cite{chaloner1984optimal,chaloner1995bayesian}．
ラベルを使わない選び方でも，たとえば点を散らすことで事後分散を減らすことはできるが，$E_0$の向きを知らなければ，修正される部分空間を$E_0$にそろえることはできない．
命題\ref{prop:design}はこの近似に依拠するので，直接検証する．

感情ごとに$K=5$本のラベルが張る空間が捉える$\|E_0\|^2$の割合を，ラベルを使わない7つの選択規則と，すべての候補のラベルを知る教示の規則（teaching rule）の下で測った．
7つの規則は，別のランダムな順序，分類器のマージン，予測クラスの上のherding，k-centerのコアセット，典型性によるクラスタリング，白色化した座標での2つの計画である．
ランダムな選択に対する比（捉える割合をそれぞれターゲットで平均してから取った比）で測ると，ラベルを使わない規則が捉える量は\NUM{p2-e58-ratio-unlab-lo}--\NUM{p2-e58-ratio-unlab-hi}倍，教示の規則が捉える量は\NUM{p2-e58-ratio-lll-lo}--\NUM{p2-e58-ratio-lll-hi}倍である（図\ref{fig:capture}）．
ターゲットのラベルを使わずに，ソース間の食い違いや仮のターゲットで学んだ修正から作った部分空間は，$\|E_0\|^2$の\NUM{p2-e58-r-uaub-lo}--\NUM{p2-e58-r-uaub-hi}\%を捉える．
選択の結果もこれと一致する．
5つの値を選んだ生徒では，標準的な能動学習の規則を含め，ラベルを使わないどの規則も$K=5$でランダムな選択を1ポイント上回らない（最大の差は\NUM{p2-cond2-e39-maxgain-k5}）．
一方，より弱い生徒を5つの値を選んだ生徒に置き換えると，ランダムな選択そのものが\NUM{p2-cond2-student-random-k5-lo}--\NUM{p2-cond2-student-random-k5-hi}ポイント上がる．
したがって，よい出発点に残る向きの誤差はターゲットに固有であり，ラベルを使わない選択からは見えない．
その修正にはラベルの情報が必要であり，しかも$4K$本のラベルが与えるよりずっと多くの情報が必要である．

\begin{figure}[t]
\centering
\FIG{capture}{\columnwidth}
\caption{ラベルを使わない選択は向きの誤差を見つけない．感情ごとに$K=5$本のラベルが張る空間が捉える出発点の向きの誤差$E_0$の割合を，ランダムな選択に対する比で，ラベルを使わない7つの規則と，すべての候補のラベルを知る教示の規則（上側の参照）について，A，V，Tで示す．点はターゲットの平均どうしの比であり，ひげはターゲットごとの比の範囲を示す．10回の反復．}
\label{fig:capture}
\end{figure}

\subsection{プール全体への弱いラベル}
\label{sec:part2-weak}

プールのすべての発話を予測する先生は，数千の発話についてラベルの情報を与える．
$Q \in [0,1]^{N_u \times 4}$を，プールの$N_u$個のラベルのない発話に対する先生のクラス確率を並べた行列とする．
$Q$の各列を，どの感情も同じ合計の重みを受けるように拡大縮小し，行列全体を合計の重みが$m$になるように拡大縮小する．
生徒は，式\eqref{eq:student}にプールにわたる$\sum_{i}\sum_{c} Q_{ic}\,\ell(W\mathbf{z}_i + b, c)$を加えたものを最小化する．
先生の誤りはプールにわたって平均されて打ち消し合い，生徒を動かすのは，先生と出発点の判定が食い違う発話の集合である．
先生と$m$は，弱いラベルを使わない選択肢を含む小さな格子から，$K$ごとに仮のターゲットで選ぶ．

先生の有用さには自然な尺度がある．
$J$を，プールでの先生のUARから出発点のUARを引いた値とする．
定義により，$J$は，先生が正しく出発点が誤って分類する発話の割合から，その逆が成り立つ発話の割合を引いたものを，クラスで平均した値に等しい．
出発点にアンカーしたまま先生のラベルに合わせる生徒は，先生の優位をせいぜい近似的に再現できるだけである．そこで本論文は，定理によらない見込みとして，生徒の利得のうち先生の情報から来る部分は$J$を下回ると考える．
列の正規化は，ラベルを使わない第2の効果を加える．
出発点自身の予測を感情の間で釣り合わせ直したものが，それだけで有用な先生として働く．

同じモダリティの先生がもたらすものは少ない．
ソースごとの分類器の投票を，仮のターゲットでのそれぞれの混同の傾向で重みづけして組み合わせると信念（belief）が得られ，5つの値を選んだ出発点に対するその$J$は，WavLMの音声で\NUM{p2-e52-q-minus-p-a}ポイントである．
この信念からの弱いラベルは，5つの値を選んだ出発点を，音声では$K = 0, 1, 5$で\NUM{p2-e59k-a-k0}，\NUM{p2-e59k-a-k1}，\NUM{p2-e59k-a-k5}ポイント，顔では\NUM{p2-e59k-v-k0}，\NUM{p2-e59k-v-k1}，\NUM{p2-e59k-v-k5}ポイント上げ（10回の反復），利得は$K=25$までに消える．
音声の$K=0$では，出発点自身の予測の釣り合わせ直しが\NUM{p2-e59k-ctrl-a-k0}ポイントを占め，信念の情報は$K \le 5$にわたって\NUM{p2-e59k-info-a-lo}--\NUM{p2-e59k-info-a-hi}ポイントを加える．
これは信念の$J$を下回る．
テキストでは信念は出発点よりよくなく，利得はほぼ0である．

\subsection{別のモダリティの先生}
\label{sec:part2-crossmodal}

ターゲットの発話が他のモダリティも伴うとき，それらのモダリティを見る，ソースで学習した融合の分類器は，別の向きを知っている．
これは，タスクがラベルを共有し入力のモダリティが異なる，特権情報の設定である\cite{lopezpaz2016unifying,vapnik2009new}．
単一モダリティの生徒に対しては，その構成の融合の出発点を先生とする．
この先生は\ref{sec:part2-start}節の手順でソースだけから学習し，ターゲットのデータで学び直さない．
配備した生徒は自分のモダリティだけを使う．
この先生の$J$は大きい．
AVの構成での2回の反復による予備的な確認では，プールで融合の出発点は音声の出発点をUARで\NUM{p2-e60pre-juar-av-a-fused}ポイント上回り，モダリティごとの予測の他の組み合わせは音声の出発点を\NUM{p2-e60pre-juar-av-a-lo}--\NUM{p2-e60pre-juar-av-a-hi}ポイント上回る．
一方，顔の出発点だけでは音声の出発点を\NUM{p2-e60pre-juar-av-a-v}ポイント下回る．
融合の生徒に対しては，単一モダリティの予測の組み合わせは融合の出発点とほとんど変わらない（同じ確認で，AVでの$J$は\NUM{p2-e60pre-juar-av-f-poe}から\NUM{p2-e60pre-juar-av-f-pool}ポイント）．
そこで2段階の参照（two-stage reference）を用いる．
まず単一モダリティの生徒をターゲットのプールで融合の先生から学習し，次にそれらの予測の積で融合の生徒を教える．

$K \le 5$では，別のモダリティからの弱いラベルは，すべての生徒を5つの値を選んだ出発点より上に上げる（図\ref{fig:ladder}に音声の生徒を示す）．
$K=0$での単一モダリティの生徒の利得は，AVで\NUM{p2-e60-av-a-k0}（音声）と\NUM{p2-e60-av-v-k0}ポイント（顔），VTで\NUM{p2-e60-vt-t-k0}（テキスト），ATで\NUM{p2-e60-at-t-k0}（テキスト），AVTで\NUM{p2-e60-avt-a-k0}（音声），\NUM{p2-e60-avt-t-k0}（テキスト），\NUM{p2-e60-avt-v-k0}ポイント（顔）である．
\NUM{p2-e60-m1-n}の単一モダリティの生徒はすべて，$K=0$と$K=5$の両方で利得を得る．
融合の生徒は$K=0$で\NUM{p2-e60-f-k0-lo}--\NUM{p2-e60-f-k0-hi}ポイントの利得を得る．
音声とテキストの生徒では，これらの利得は同じモダリティの先生による利得を上回る．
顔の生徒では小さい（$K=0$で，VTとAVTでの\NUM{p2-e60-vt-v-k0}ポイントと\NUM{p2-e60-avt-v-k0}ポイントに対し，同じモダリティの先生では\NUM{p2-e59k-v-k0}ポイント）．
$J$を測ったAVでは，前の小節の尺度が見込むとおり，音声の生徒の利得はその先生の$J$を下回る．
ラベルが増えると利得は縮み，$K \ge 10$では\NUM{p2-e60-m3-students-n}の生徒のうち\NUM{p2-e60-m3-loss-n}が5つの値を選んだ出発点を下回る（最も悪い差は\NUM{p2-e60-m3-worst}ポイント）．
拡張した候補群（extended family）は，先生の候補として$K$本のラベルで学び直した融合の生徒と2段階の参照を加え，弱いラベルの重みづけ（釣り合わせる，釣り合わせない，ハード）を選択の対象に加えたものである．
拡張した候補群は，比較で評価するAVTの音声の生徒とAVの融合の生徒について，これらの損失をなくす（\ref{sec:comparison}節）．
VTのテキストの生徒とATの音声の生徒は，拡張した候補群では実行していない．

\begin{figure}[t]
\centering
\FIG{ladder}{\columnwidth}
\caption{音声の生徒における提案法の段階ごとの結果．(a)~ターゲットのみの学習，合併（単一のソースにわたる平均），合併したソースからの仮説転移（HT），5つの値を選んだ出発点のUARを$K$に対して示す（7つのターゲット，30回の反復）．(b)~同じ実験の5つの値を選んだ出発点からの差を，同じモダリティの信念からの弱いラベル（7つのターゲット），AVの融合の先生（6つのターゲット）とAVTの融合の先生（4つのターゲット）からの弱いラベル，AVTでの拡張した候補群（4つのターゲット）について，10回の反復で示す．融合の構成はそのすべてのモダリティがそろった発話だけを残すので，出発点は構成ごとに異なる．}
\label{fig:ladder}
\end{figure}
